\documentclass[varwidth=14.5cm,border=10pt]{standalone}
\usepackage{amsmath,xcolor,array}
\usepackage{fontspec,unicode-math}
\setmainfont{DejaVu Sans}[ItalicFont={DejaVu Sans},ItalicFeatures={FakeSlant=0.2},BoldItalicFont={DejaVu Sans Bold},BoldItalicFeatures={FakeSlant=0.2}]
\setmathfont{texgyredejavu-math.otf}
\setmathfont{DejaVu Sans}[range={up,bfup}]
\setmathfont{DejaVu Sans}[range={it,bfit},FakeSlant=0.2]
\setmathfont{texgyredejavu-math.otf}[range={"2202,"1D715,"03BB,"03B4,"03B8,"2207,"2211}]
\pagecolor[HTML]{1E1E2E}\color[HTML]{F4F4F4}
\definecolor{accent}{HTML}{89B4FA}
\definecolor{warn}{HTML}{F38BA8}
\definecolor{good}{HTML}{A6E3A1}
\definecolor{ref}{HTML}{9399B2}
\newcommand{\heading}[1]{\par\medskip{\large\bfseries\color{accent}#1}\par\smallskip}
\newcommand{\src}[1]{{\color{ref}\small[#1]}}
\setlength{\parindent}{0pt}
\setlength{\tabcolsep}{5pt}
\raggedright
\begin{document}

{\Large\bfseries The forward pass and the backward pass}

\heading{The one idea}
They are the \textbf{same chain of operations, walked in opposite directions.}
\[
\begin{array}{ccccccccc}
 x &\to& a_j &\to& h_j &\to& \hat y &\to& J\\
 \multicolumn{9}{c}{\color{good}\text{\textbf{FORWARD}: carries VALUES --- what the network actually computes}}\\[10pt]
 \multicolumn{9}{c}{\color{warn}\text{\textbf{BACKWARD}: carries SENSITIVITIES --- how much } J \text{ cares}}\\
 \delta_x &\leftarrow& \delta_{a_j} &\leftarrow& \delta_{h_j} &\leftarrow& \delta_{\hat y} &\leftarrow& 1
\end{array}
\]
Forward you ask \emph{"what comes out?"}. Backward you ask \emph{"what would change it?"}

\heading{Forward --- left to right, one evaluation}
Plug in $x$ and push. Nothing here is a derivative; it is just the model being used.
\[
\begin{array}{l@{\ \ }l@{\qquad}l}
 a_j = w_j x + b_j & \text{weigh and offset} & (+0.500,\ -0.250,\ -0.375)\\
 h_j = \tanh a_j & \text{bend} & (+0.462117,\ -0.244919,\ -0.358357)\\
 \hat y = \sum_j v_j h_j + c & \text{recombine} & +0.389291\\
 r = \hat y - y & \text{the miss} & +1.389291\\
 J = r^2 & \text{the loss} & 1.930130
\end{array}
\]
{\color{good}This is exactly what you do when you \emph{use} a trained network}
(\texttt{yq = net(xQuery')}). Training just does it with the loss attached at the end.\\
{\color{ref}Store $h_j$ on the way --- the backward pass needs it.}

\heading{Backward --- right to left, one sweep}
Start at the end with $\partial J/\partial J=1$, and at each step \textbf{multiply by that
link's local derivative}. Write $\delta_{(\cdot)}=\partial J/\partial(\cdot)$.
\[
\begin{array}{r@{\ }c@{\ }l@{\qquad}l}
 \delta_{\hat y} &=& 2r & = +2.778583\\[4pt]
 \delta_{h_j} &=& \delta_{\hat y}\cdot v_j & = (+1.389291,\ -2.778583,\ +2.222866)\\[4pt]
 \delta_{a_j} &=& \delta_{h_j}\cdot\big(1-h_j^2\big) & = (+1.092605,\ -2.611909,\ +1.937406)
\end{array}
\]
\[
  \text{\color{ref}using } 1-h_j^2 = (0.786448,\ 0.940015,\ 0.871580)
\]
Each $\delta$ answers: \emph{if this intermediate number were nudged by 1, how much would
$J$ move?} Notice $\delta_{h_2}$ is the biggest in size --- neuron 2 is the one $J$ cares
about most at this data point.

\heading{Harvesting the gradient}
A parameter sits \textbf{on} a link. Its derivative is the $\delta$ arriving there times its
own local factor --- no further chaining:
\[
\begin{array}{l@{\qquad}l@{\qquad}l}
 \dfrac{\partial J}{\partial v_j}=\delta_{\hat y}\,h_j &
 \dfrac{\partial J}{\partial c}=\delta_{\hat y} &
 \text{\color{ref}harvested at }\hat y\\[10pt]
 \dfrac{\partial J}{\partial w_j}=\delta_{a_j}\,x &
 \dfrac{\partial J}{\partial b_j}=\delta_{a_j} &
 \text{\color{ref}harvested at }a_j
\end{array}
\]
So $\partial J/\partial w_1 = 1.092605\times 0.75 = +0.819454$, and the ten derivatives
drop out as the sweep passes.

\heading{Why forward must come first}
Every local derivative in the backward pass is {\color{warn}evaluated at the forward
values}: $\partial\hat y/\partial h_j = v_j$ needs nothing, but
$\partial h_j/\partial a_j = 1-h_j^2$ needs $h_j$, and
$\partial J/\partial\hat y = 2r$ needs $r$.
No forward pass, no numbers to differentiate \emph{at}. That is also the memory cost of
training: you must keep the whole forward trace alive until the backward sweep has consumed it.

\heading{You already ran both in A3}
\[
\begin{array}{l@{\qquad}l@{\qquad}l}
 & \textbf{truss (L06)} & \textbf{network (L07)}\\[3pt]
 \text{forward} & \text{solve } K(a)q=f \to q & \text{evaluate } x\to\hat y\\
 \text{the response} & f = f^{\!\top}q\ \text{(compliance)} & J = r^2\ \text{(loss)}\\
 \text{backward} & \text{adjoint } \lambda\ (=-q\ \text{here}) & \text{the } \delta\text{'s}\\
 \text{harvest} & \partial f/\partial a_e = -q_e^{\!\top}K_e^{\text{unit}}q_e
   & \partial J/\partial\theta_k=\delta\cdot(\text{local})\\[5pt]
 \text{cost} & \text{2 solves, any } n_e & \text{2 sweeps, any } |\theta|
\end{array}
\]
{\color{good}Same two-pass shape. One evaluation of the physics, one sweep back against it,
then every derivative falls out for free.}

\end{document}
